% !TeX root = ../navier-stokes-manuscript.tex
\subsection{The case we are dealing with}\label{sub:ThePerfectlyStretchingVortex}

The case we are dealing with is a vortex that stays aligned with extensional strain and keeps stretching.
Suppose its vorticity remains along a positive extensional direction of the strain, so the strain amplifies it while incompressibility contracts the cross-section of the material filament.
The vortex narrows, its gradients sharpen, and viscosity acts on those same gradients.
Bending, wobble, loss of alignment, and fragmentation would interrupt the motion.
Set them aside.
The hard case is the vortex that stays in the stretch.

The question is immediate: does moving this profile to a smaller scale change the competition between stretching and viscosity?
The Navier--Stokes scaling gives them exactly the same relative weight.

Here is the balance itself.
Write
\[
  \omega:=\nabla\times u,
  \qquad
  S:=\frac12\bigl(\nabla u+(\nabla u)^{\mathsf T}\bigr).
\]
Taking the curl of Navier--Stokes gives
\begin{equation}\label{eq:VorticityEquation}
  \partial_t\omega+(u\cdot\nabla)\omega
  =S\omega+\nu\Delta\omega.
\end{equation}
When \(\omega\) points along an eigenvector of \(S\) with positive eigenvalue, \(S\omega\) amplifies it.
For an aligned material filament, that extension contracts the cross-sectional area and concentrates the vorticity.
The term \(\nu\Delta\omega\) diffuses the gradients created by the same concentration.

Now shrink the full profile.
To ask what scale alone does to the contest, replace a characteristic length \(\ell\) by \(\ell/\lambda\), with \(\lambda>1\), through the Navier--Stokes rescaling
\[
  u_\lambda(x,t)=\lambda u(\lambda x,\lambda^2t),
  \qquad
  p_\lambda(x,t)=\lambda^2p(\lambda x,\lambda^2t).
\]
Every term in the momentum equation then acquires the same factor:
\[
\begin{aligned}
  \partial_tu_\lambda
  &=\lambda^3(\partial_tu)(\lambda x,\lambda^2t),
  &
  (u_\lambda\cdot\nabla)u_\lambda
  &=\lambda^3((u\cdot\nabla)u)(\lambda x,\lambda^2t),
  \\
  \nabla p_\lambda
  &=\lambda^3(\nabla p)(\lambda x,\lambda^2t),
  &
  \nu\Delta u_\lambda
  &=\lambda^3\nu(\Delta u)(\lambda x,\lambda^2t).
\end{aligned}
\]
The vortex equation says the same thing in the variables of the physical picture:
\[
  \omega_\lambda(x,t)=\lambda^2\omega(\lambda x,\lambda^2t),
  \qquad
  S_\lambda(x,t)=\lambda^2S(\lambda x,\lambda^2t).
\]
Both \(S_\lambda\omega_\lambda\) and \(\nu\Delta\omega_\lambda\) carry the factor \(\lambda^4\).
Going to a smaller scale preserves their relative weight.
That is the critical balance.

To turn that balance into a height, we need a norm that assigns the same size to the original profile and its concentrated rescaling.
Since
\[
  \widehat{u_\lambda}(\xi,t)
  =\lambda^{-2}\widehat u(\xi/\lambda,\lambda^2t),
\]
the change of variables \(\eta=\xi/\lambda\) gives
\begin{align*}
  \norm{u_\lambda(t)}_{\dot H^s}^2
  &=
  \int_{\R^3}
  |\xi|^{2s}\lambda^{-4}
  \left|\widehat u(\xi/\lambda,\lambda^2t)\right|^2
  \,d\xi
  \\
  &=
  \lambda^{2s-1}
  \int_{\R^3}
  |\eta|^{2s}
  \left|\widehat u(\eta,\lambda^2t)\right|^2
  \,d\eta.
\end{align*}
Hence
\[
  \norm{u_\lambda(t)}_{\dot H^s}
  =\lambda^{s-1/2}\norm{u(\lambda^2t)}_{\dot H^s}.
\]
There is exactly one exponent at which the scale factor disappears: \(s=\tfrac12\).
With \(\Lambda:=(-\Delta)^{1/2}\), set
\begin{equation}\label{eq:BodyCriticalHeightDefinition}
  H(t)
  :=\frac12\norm{u(t)}_{\dot H^{1/2}}^2
  =\frac12\norm{\Lambda^{1/2}u(t)}_2^2.
\end{equation}
This is the critical height.
The profile can move to a shorter length while keeping the same value of \(H\).

Now put kinetic energy beside it.
At \(s=0\),
\[
  \norm{u_\lambda(t)}_2^2
  =\lambda^{-1}\norm{u(\lambda^2t)}_2^2.
\]
The energy identity comes from pairing Navier--Stokes with \(u\): incompressibility removes the transport and pressure pairings, while integration by parts gives the viscous sign.
Thus
\[
  \frac12\frac{d}{dt}\norm{u(t)}_2^2
  +\nu\norm{\nabla u(t)}_2^2
  =0,
\]
and after integration in time,
\begin{equation}\label{eq:BodyEnergyIdentity}
  \frac12\norm{u(t)}_2^2
  +\nu\int_0^t\norm{\nabla u(s)}_2^2\,ds
  =\frac12\norm{u_0}_2^2.
\end{equation}
The law controls total kinetic energy.
Under the concentration we are testing, that energy is multiplied by \(\lambda^{-1}\) while the critical height is unchanged.
The bound gets easier to satisfy as the vortex narrows.
This is why energy leaves the regularity problem open: it controls a quantity that becomes smaller along the very concentration that \(H\) continues to see.

The scale can shrink while remaining inside the energy bound.
What matters next is the time available at that scale.
A structure of length \(\ell\) has the linear viscous smoothing time
\begin{equation}\label{eq:BodyHeatComparisonTime}
  \delta t_{\mathrm{heat}}(\ell)\asymp\frac{\ell^2}{\nu}.
\end{equation}
Replacing \(\ell\) by \(\ell/\lambda\) divides this time by \(\lambda^2\), exactly the time compression in \(u_\lambda(x,t)\).

Take a geometric scale family
\[
  \ell_k=\rho^k\ell_0,
  \qquad 0<\rho<1.
\]
With scale-uniform comparison constants,
\[
  \delta t_k^{\mathrm{heat}}
  \asymp\frac{\ell_0^2}{\nu}\rho^{2k},
  \qquad
  \sum_{k=0}^{\infty}\delta t_k^{\mathrm{heat}}
  \asymp
  \frac{\ell_0^2}{\nu}\frac{1}{1-\rho^2}<\infty.
\]
So far this is a comparison clock.
The solution supplies a cascade when it actually produces the smaller states on that clock.

The scale family tells us how small the active region could become; \(H\) tells us whether the evolving solution is actually escaping through critical height.
Mark that escape by first passages.
Suppose \(H\) crosses arbitrarily large levels, choose \(H_0>H(0)\), and define
\[
  L_k:=2^kH_0,
  \qquad
  t_k:=\inf\{t>0:H(t)=L_k\}.
\]
Continuity gives \(t_0<t_1<t_2<\cdots\).
Suppose also that the \(k\)-th passage is carried by an active length \(\ell_k=\rho^k\ell_0\) and obeys the uniform comparison
\[
  0<t_{k+1}-t_k
  \leq C\,\delta t_k^{\mathrm{heat}}
  \asymp C\frac{\ell_0^2}{\nu}\rho^{2k}.
\]
Then
\[
  \sum_{k=0}^{\infty}(t_{k+1}-t_k)<\infty,
  \qquad
  t_k\uparrow\tau<\infty.
\]
Now the comparison clock has become the Zeno route in the evolving solution: the critical height crosses every chosen level while the active scale and the time between passages both collapse.

The same route sharpens the spatial profile.
For every \(m\in\Nzero\),
\[
  \nabla^m u_\lambda(x,t)
  =\lambda^{m+1}(\nabla^m u)(\lambda x,\lambda^2t).
\]
At the \(k\)-th scale, let \(\lambda_k:=\ell_0/\ell_k=\rho^{-k}\).
The gradient carries the factor \(\lambda_k^2\).
If the realized states retain a nonzero scale-normalized gradient profile, their gradients diverge while the passage times converge to \(\tau\).
That limiting gradient failure is the singularity direction of the perfectly stretching vortex.

The field still has to make every one of those passages.
A shorter interval magnifies the temporal response whenever a definite change survives across it.
This is where velocity, acceleration, jerk, and the later responses enter.
